Here, we discuss the general scientific objectives of the project. Ultimately, as was discussed before, we want to be able to express the asymptotic of the partition function for some model of the external field ensemble with as much precision as possible. Let's lay out the path of this project to obtain such a goal.

Our main object is the partition function $\mathcal{Z}$ and at the same time, the free energy $\ln(\mathcal{Z})$. If we think of models described so far, we can see that generally it is possible to express such a function in terms of the associated norming constants, i.e., in terms of the orthogonal polynomial associated with the problem and consequently, because of the reproducing property, the norming constant. For this work we will deal with what was called the bi-orthogonal ensemble. Naturally because of the different mathematical structure, the partition function won't have the same expression as before in terms of norming constants. It is expected however to be some analogous formula. Our first goal is then to find an expression of the partition function in terms of these norming constants for the bi-orthogonal polynomials that define the problem. 

\begin{tcolorbox}
	Express the partition function in terms of the norming constants for the bi-orthogonal polynomials that define the problem of the bi-orthogonal ensemble
\end{tcolorbox}

Interestingly, such polynomials have an accessible asymptotic. This was explored and determined in a recent work \cite{Bleher_2006}; Pavel and Arno have found the asymptotic for the large $n$ limit of Gaussian random matrices with external source, in terms of the associated Riemann-Hilbert problem for the model. So, for this step, the objective is to translate the results into the asymptotic of the norming constants and plug it into our know partition function formula developed in \cite{Byun_2023}. Finally, possessing the asymptotic for the bi-orthogonal polynomials we should be able with the aid of standard techniques to express the asymptotic for $\mathcal{Z}$.  This is therefore a very natural final objective.

\begin{tcolorbox}
	Connect the asymptotic given by the Riemann-Hilbert problem by Pavel and Arno into the asymptotic for the bi-orthogonal polynomials using the formula in \cite{Byun_2023} to get out asymptotic for the partition function.
\end{tcolorbox}

Possessing such asymptotic we can validate and check what kind of gains or loses we had in applying the method described in \cite{Byun_2023} and studied in the student's masters. Therefore checking for the usefulness of such technique in this relatively new model.